%%%%%%%%%%%%%%%%%%%%%%%%%%%%%%%%%%%%%%%%%
% imec System Software Engineer cover letter - Arman Babaei
%%%%%%%%%%%%%%%%%%%%%%%%%%%%%%%%%%%%%%%%%
\documentclass[11pt]{resume}
\usepackage{ebgaramond}
\usepackage[normalem]{ulem}
\usepackage[hidelinks]{hyperref}
\AtBeginDocument{\let\oldhref=\href\renewcommand\href[2]{\dotuline{\oldhref{#1}{#2}}}}
\hypersetup{pdftitle={Arman Babaei - imec Cover Letter},pdfauthor={Arman Babaei}}
\name{Arman Babaei}
\address{+98~$\cdot$~919~$\cdot$~043~$\cdot$~7966 \\ 292.arma@gmail.com \\ \href{https://www.linkedin.com/in/arman-babaei/}{linkedin} \\ \href{https://github.com/Arman17Babaei}{github}}

\begin{document}
\vspace{1em}
Hiring Team\\
imec\\
Leuven, Belgium

\textbf{Application for System Software Engineer}

Dear Hiring Team,

I am applying for the System Software Engineer position in Leuven. My experience combines professional C++ and Python development with experimental systems research at Sharif University of Technology, EPFL and ISTA. I am interested in the practical work at the center of this role: implementing software that makes a research idea testable and integrating it into a working system.

At Mahsan, I developed parallel network-relaying software in C++ that connected to ISP switches. I also led development of a Python domain-specific language for end-to-end testing using simulated network packets, and served as Scrum Master for a team of five. More recently, at Tapsi, I have worked across backend development, site reliability and platform engineering. This has included rewriting production services in Go, investigating failures across services and networking, and developing deployment automation with canary releases and rollback controls.

My master's research focuses on low-latency networking for volumetric video streaming. I implemented a DualPI2 prototype in P4/BMv2 and compared it with a Linux reference implementation. I also built a reproducible test harness combining C traffic fixtures, Python analysis and QEMU/Mininet automation. The work connects protocol behavior and packet-level measurements to application delivery and rendered quality, requiring careful validation across several interacting components. My thesis has been approved, and the oral defense is pending scheduling.

At EPFL, I extended a Go simulator to study fairness in decentralized storage. At ISTA, I built an experimental test bench and contributed to work on parallel smart-contract execution. These collaborations gave me experience translating research questions into implementations and measurable experiments. My project work also includes contributions to a C-based UEFI system information viewer, specifically ACPI inspection and network-interface discovery through firmware protocols.

I would bring imec experience in systems programming, test automation and production debugging, together with an interest in learning the hardware and experimental context of each project. The opportunity to work with software, electronics and system engineers on research prototypes is particularly appealing to me.

I am available to start from November 2026, subject to obtaining the necessary visa and work authorization for Belgium, for which I would require employer support. I would welcome a conversation about how my experience could contribute to your team.

Kind regards,\\[0.4em]
Arman Babaei
\end{document}
